\scanpage{42}
\mspage{\unsure{1}}

% B01 | head, lines 1–4
Soit $B$ un anneau avec unité et un \unsure{gpe|groupe} fini $G$ d'automorphismes $A=B^G$. On pose $\operatorname{Tr}_G b=\sum_{g\in G}gb$. Conditions équivalentes, ($q$ est un entier donné) : $A$ est un sous-anneau, et $\operatorname{Tr}_G:B\to A$ est un \unsure{hom} de bi-$A$-modules

% B02 | left of opening list, read clockwise
% margin note written sideways
\margin{\unsure{Voy. dans le Commentaire|Voy. dans le} \unsure{précis ou corollaire Voy. de Serre|précis au corollaire Voy. de Serre} Comment voir ça}

% B03 | four conditions at upper left
\[
\begin{aligned}
\text{a)}\quad & A=\operatorname{Tr}_G B\qquad\text{i.e. }\widehat H^0(G,B)=\{0\}\\
\text{a$'$)}\quad & 1=\operatorname{Tr}_G b\qquad\text{avec }b\in B\\
\text{b)}\quad & \widehat H(G,M)=0\qquad\text{pour }\unsure{\forall|\text{tout}}\ (B,G)\text{-module }M\\
\unsure{b(q)|b_q}\text{)}\quad & \widehat H^q(G,M)=0
\end{aligned}
\]

% B04 | Démonstration, first four lines
Démonstration : $a)\Longrightarrow a')$ trivial, $a'\Longrightarrow b$ car on a, en désignant par $\varphi_M(b)$ l'homothétie dans $M$ définie par un $b\in B$ : $g\varphi_M(b)g^{-1}=\varphi_M(gb)$ d'où $\varphi_M(\operatorname{Tr}_G b)=\operatorname{Tr}_G\varphi_M(b)$, d'où le

% B05 | first Corollaire and bracketed qualification
\uline{Corollaire} Si $\lambda\in\operatorname{Tr}_G B$, l'homothétie $\lambda$ dans $\widehat H(G,M)$ est \uline{nulle} \inserted{[et la condition est \unsure{manifeste\textsuperscript{t}|manifestement} nécessaire, en prenant $M=\underline B$)} qui implique $a'\Longrightarrow b$. $b\Longrightarrow b(q)$ trivial. Or $b(q)\Longrightarrow b(q+1)$

% B06 | immediately before sequence
car on a

% B07 | at fold, exact sequence and assignment below
\[
\begin{aligned}
0\to M&\to M\otimes_Z Z(G)\to N\to 0\\[-2pt]
m&\to\textstyle\sum m\otimes g
\end{aligned}
\]

% B08 | four lines under sequence
où on met sur le terme milieu la \deleted{\unsure{$B$-}}structure de $B$-module définie par $\lambda(m\otimes g)=\lambda m\otimes g$, et la structure de $G$-module définie par $g'(m\otimes g)=g'm\otimes g'g$, alors $g'.(\lambda[m\otimes g])=(g'.\lambda)(g'.[m\otimes g])$ car les deux mbres sont \unsure{respt} $g'.(\lambda m\otimes g)=g'(\lambda m)\otimes g'g$ et $(g'\lambda)(g'm\otimes g'g)=(g'\lambda)(g'm)\otimes g'\unsure{g|m}$, c'est la $=$ chose car $g'.(\lambda m)=(g'.\lambda)(g'.m)$.

% B09 | dimension shift, down to second corollary
Alors on a, puisque $\widehat H(G,M\otimes_Z Z(G))=0$, $\widehat H^{q+1}(G,M)=\widehat H^q(G,N)=0$. On prouve de façon analogue $b(q)\Longrightarrow b(q-1)$, d'où $b(q)\Longrightarrow b(0)$. Mais $b(0)\Longrightarrow a$ trivialement.

% B10 | second Corollaire through bottom
\uline{Corollaire} Supposons $B$ \inserted{commutatif} intègre et intégralement clos, \inserted{$A$ local.} Les conditions de la proposition sont vérifiées si l'indice de ramification $e$ de $B$ sur $A$ ($={}$ ordre du groupe d'inertie) \deleted{\unsure{soit premier à la|est premier à la} \inserted{\unsure{divisible}} \unsure{premier à}} \deleted{la caractéristique \unsure{(supposée|résiduelle} $\ne0$) \unsure{par} revient à écrire un $u\in B$ tel que $\unsure{eu=1}$}. est inversible dans $A$, i.e. \unsure{non nul|une unité} dans $A/\unsure{\mathfrak m(A)}$ \deleted{\illegible{1}}

% archive sheet number: 41

